% Generated draft: requires mathematical and source review.
\documentclass[11pt]{article}
\usepackage[margin=25mm]{geometry}
\usepackage{fontspec}
\setmainfont{Latin Modern Roman}
\usepackage{amsmath,amssymb,mathtools,unicode-math}
\setmathfont{Latin Modern Math}
\usepackage{enumitem,needspace,xurl,hyperref}
\setlength{\parindent}{0pt}
\setlength{\parskip}{5pt}
\setlength{\emergencystretch}{3em}
\newcommand{\N}{\mathbb{N}}
\newcommand{\Z}{\mathbb{Z}}
\newcommand{\Q}{\mathbb{Q}}
\newcommand{\R}{\mathbb{R}}
\newcommand{\C}{\mathbb{C}}
\newcommand{\F}{\mathbb{F}}
\newcommand{\A}{\mathbb{A}}
\newcommand{\PP}{\mathbb{P}}
\newcommand{\E}{\mathbb{E}}
\newcommand{\eps}{\varepsilon}
\newcommand{\dd}{\,\mathrm{d}}
\begin{document}
\section*{UnsolvedMath ID: 30006556}
\textbf{Source problem code:} catalogue entry
\subsection{Definitions and mathematical statement}
Let \(n\in\N\) with \(n>4\), and let \(P\subset\R^2\) be a set of \(n\) distinct points, with Euclidean distance \(\|p-q\|_2\). Define the set of distances determined by \(P\) by
\[
D(P)=\{\|p-q\|_2:p,q\in P,\ p\ne q\}.
\]
For each \(d\in D(P)\), define its multiplicity by
\[
m_P(d)=\#\bigl\{\{p,q\}:p,q\in P,\ p\ne q,\ \|p-q\|_2=d\bigr\}.
\]
Thus each unordered pair of distinct points is counted once, and every distance in \(D(P)\) has positive multiplicity.

\textbf{Question.} Is it true that for every integer \(n>4\) and every such set \(P\), there exist \(d_1,d_2\in D(P)\) such that
\[
d_1\ne d_2,\qquad 1\le m_P(d_1)\le n,\qquad 1\le m_P(d_2)\le n?
\]
The question imposes no convexity or general-position assumption on \(P\). The supplied background mentions results for convex or ``not-too-convex'' sets, but does not define the latter class or state the precise hypotheses of those results; those qualifications cannot be formalized further from this record.
\subsection{Short English statement}
Must every set of \(n>4\) distinct points in the Euclidean plane determine two different positive distances, each realized by at most \(n\) unordered pairs of points? The supplied record labels the general question open; this status is not independently verified here.
\subsection{Sources}
\begin{itemize}
\item Supplied local UnsolvedMath registry record, permanent ID \(30006556\), title: \emph{Two Low-Multiplicity Distances in Planar Point Sets}; source problem code \texttt{catalogue entry}, in the Oberwolfach Reports Open Problems collection. The record gives proposed year 2026 but names no proposer.
\item \emph{Combinatorics} (2026), Oberwolfach Reports, Report 1/2026, pp.~5--79, DOI: \url{https://doi.org/10.4171/owr/2026/1}. This is the supplied report attribution; no exact page for this question is supplied. The registry labels the statement \texttt{corrected\_verified} and reports that the official report was checked, the question isolated from surrounding exposition, and its notation repaired. These are supplied verification claims, not checks performed for this draft.
\item Clemen--Dumitrescu--Liu, \emph{On the multiplicities of interpoint distances}, \url{https://arxiv.org/abs/2505.04283}.
\item Supplied literature assessment dated exactly \texttt{2026-08-22}: the distance-multiplicity paper is reported to establish the claim for convex or ``not-too-convex'' sets, while the universal assertion for arbitrary planar point sets is reported to remain open. The supplied sources checked are the Oberwolfach report and the Clemen--Dumitrescu--Liu paper linked above. This is supplied dated curation derived from the OWR record, not an independently verified literature assessment; time-sensitive status and the cited special-case claims require independent verification.
\end{itemize}
\end{document}
